% Options for packages loaded elsewhere
\PassOptionsToPackage{unicode}{hyperref}
\PassOptionsToPackage{hyphens}{url}
%
\documentclass[
  10pt,
]{article}
\usepackage{amsmath,amssymb}
\usepackage{iftex}
\ifPDFTeX
  \usepackage[T1]{fontenc}
  \usepackage[utf8]{inputenc}
  \usepackage{textcomp} % provide euro and other symbols
\else % if luatex or xetex
  \usepackage{unicode-math} % this also loads fontspec
  \defaultfontfeatures{Scale=MatchLowercase}
  \defaultfontfeatures[\rmfamily]{Ligatures=TeX,Scale=1}
\fi
\usepackage{lmodern}
\ifPDFTeX\else
  % xetex/luatex font selection
\fi
% Use upquote if available, for straight quotes in verbatim environments
\IfFileExists{upquote.sty}{\usepackage{upquote}}{}
\IfFileExists{microtype.sty}{% use microtype if available
  \usepackage[]{microtype}
  \UseMicrotypeSet[protrusion]{basicmath} % disable protrusion for tt fonts
}{}
\makeatletter
\@ifundefined{KOMAClassName}{% if non-KOMA class
  \IfFileExists{parskip.sty}{%
    \usepackage{parskip}
  }{% else
    \setlength{\parindent}{0pt}
    \setlength{\parskip}{6pt plus 2pt minus 1pt}}
}{% if KOMA class
  \KOMAoptions{parskip=half}}
\makeatother
\usepackage{xcolor}
\usepackage[margin=0.85in]{geometry}
\setlength{\emergencystretch}{3em} % prevent overfull lines
\providecommand{\tightlist}{%
  \setlength{\itemsep}{0pt}\setlength{\parskip}{0pt}}
\setcounter{secnumdepth}{-\maxdimen} % remove section numbering
\usepackage{amsmath,amssymb,amsthm}
\ifLuaTeX
  \usepackage{selnolig}  % disable illegal ligatures
\fi
\usepackage{bookmark}
\IfFileExists{xurl.sty}{\usepackage{xurl}}{} % add URL line breaks if available
\urlstyle{same}
\hypersetup{
  pdftitle={The image of the finite-energy Monge--Ampere operator on a polarized affine cone},
  pdfauthor={Research manuscript},
  hidelinks,
  pdfcreator={LaTeX via pandoc}}

\title{The image of the finite-energy Monge--Ampere operator on a
polarized affine cone}
\usepackage{etoolbox}
\makeatletter
\providecommand{\subtitle}[1]{% add subtitle to \maketitle
  \apptocmd{\@title}{\par {\large #1 \par}}{}{}
}
\makeatother
\subtitle{Actual realization at every fixed Reeb field}
\author{Research manuscript}
\date{5 October 2026}

\begin{document}
\maketitle

\section{1. Statement, domain, and
conventions}\label{statement-domain-and-conventions}

Let \(k\) be an arbitrary algebraically closed field of characteristic
zero, equipped with its trivial absolute value. Let
\(X=\operatorname{Spec}R\) be a normal good affine cone, with its
prescribed algebraic torus \(T\). Thus \(R\) is a finitely generated
normal domain with a character decomposition, the degree-zero invariant
algebra is \(k\), and the grading has a unique attractive vertex. A Reeb
vector is a real vector \(\xi\) in the cocharacter space such that
\(\langle\xi,\alpha\rangle>0\) for every nonzero character occurring in
\(R\). The Reeb vector in the theorem need not be rational. Put
\(n=\dim X\geq1\) and \(p=n-1\).

We retain the original invariant carrier \[
                         K=X_0^T.
\] It is a compact Hausdorff space. The superscript denotes
torification; it does not mean that we have enlarged the carrier to the
full analytic link. All the cone potentials below are the actual
equivariant potentials on the punctured analytic cone, restricted to
\(K\). Their full functions are determined by these restrictions,
equivariance, and torification.

For example, the actual Fubini--Study class \(H(\xi)\) consists of
finite primary homogeneous maxima \[
 u(v)=\max_i\frac{\lambda_i-v(f_i)}
                         {\langle\xi,\operatorname{wt}(f_i)\rangle},
 \qquad \bigcap_i\{f_i=0\}=\{o\},\quad \lambda_i\in\mathbb R.
                                                               \tag{1.1}
\] \(\operatorname{CPSH}(\xi)\) is its local uniform closure.
\(\operatorname{PSH}(\xi)\) consists of the nontrivial pointwise
decreasing net limits of these genuine continuous potentials. In
particular, a potential can take \(-\infty\); it is not an element of an
abstract metric completion. Let \(\rho_\xi\) be the original reference,
so \(\rho_\xi|_K=0\), and let \(V_\xi\) be the original cone volume. The
genuine mixed probability is \[
 \operatorname{MA}_\xi(u_1,\ldots,u_p)
 =V_\xi^{-1}dd^cu_1\cdots dd^cu_p\,dd^c\rho_\xi^+,
 \qquad \rho_\xi^+=\max(\rho_\xi,0).                              \tag{1.2}
\] Its support lies in \(K\) and its mass is one. On continuous
potentials the energy is \[
 E_\xi(u)=\frac1{nV_\xi}\sum_{j=0}^p
       \int_Ku\,(dd^cu)^j(dd^c\rho_\xi)^{p-j}dd^c\rho_\xi^+.
                                                               \tag{1.3}
\] On PSH it is the infimum of the energies of dominating actual
continuous potentials. It is increasing, concave, translation
equivariant, and continuous along decreasing nets. Set
\(\mathcal E^1(\xi)=\{u\in\operatorname{PSH}(\xi):E_\xi(u)>-\infty\}\).
Formula (1.2) on this class means the existing fixed-\(\xi\)
decreasing-net extension, including its weighted integral extension. We
use \[
 u(v_{\rm triv})=\sup_Ku=0,\qquad J_\xi(u)=-E_\xi(u)               \tag{1.4}
\] for normalization. The reference probability is
\(\operatorname{MA}_\xi(\rho_\xi)=\delta_{v_{\rm triv}}\).

For a positive Radon probability \(\mu\) on \(K\), define intrinsically
\[
 L_\xi(\mu)=\sup_{h\in H(\xi)}
       \left(E_\xi(h)-\int_Kh\,d\mu\right).                       \tag{1.5}
\] The definition uses bounded, actual potentials, and so is meaningful
without assuming an inverse equation. Constants cancel in (1.5).

\textbf{Theorem.} For every cone and every fixed Reeb vector specified
above, \[
 \operatorname{MA}_\xi\bigl(\mathcal E^1(\xi)/\mathbb R\bigr)
 =\{\mu\text{ positive Radon on }K:\mu(K)=1, L_\xi(\mu)<\infty\}.
                                                               \tag{1.6}
\] The condition on the right is equivalent to integrability of every
actual \(\mathcal E^1(\xi)\) potential. Every such measure is the
genuine fixed-\(\xi\) Monge--Ampere probability of an actually attained
normalized finite-energy potential \(w\), unique modulo constants. It is
also the attained maximizer of \[
                  E_\xi(u)-\int_Ku\,d\mu.
\] For \(p\geq1\), every normalized actual \(\mathcal E^1(\xi)\)
maximizing sequence \((u_j)\), with deficits \[
 d_j=L_\xi(\mu)-E_\xi(u_j)+\int_Ku_j\,d\mu,
\] satisfies the exact estimate \[
 I_\xi(u_j,w)\leq\frac{d_j}{2a_p},\qquad
 a_p=\frac{p}{4^{p+1}},\qquad
 I_\xi(u,v)=\int_K(u-v)\,d\bigl(\operatorname{MA}_\xi(v)
                                      -\operatorname{MA}_\xi(u)\bigr).
                                                               \tag{1.7}
\] In particular \(I_\xi(u_j,w)\to0\). If \(p=1\), (1.7) reads
\(I_\xi(u_j,w)\leq8d_j\). The dimension-one cone has a one-point
carrier; there is no use of \(a_0\).

The proof uses the established projective envelope theorem stated in
Section 2 and the existing cone calculus described below. These are
mathematical background premises of this paper. The envelope property is
verified on the actual normal quotient in Section 3, so no envelope
condition remains in the theorem. Source extraction and document
compilation are documentary checks. This reviewed edition retains the
stated mathematical background; its review and build records accompany
the paper.

\section{2. The established projective envelope theorem and its precise
use}\label{the-established-projective-envelope-theorem-and-its-precise-use}

We first state the projective background with its full domain. It is
essential to distinguish this theorem from an envelope assertion on a
resolution.

\textbf{Projective envelope lemma.} Let \(Z\) be an original reduced
irreducible geometrically unibranch projective variety over any
algebraically closed trivially valued characteristic-zero field, of
arbitrary cardinality. Let \(H\) be an actual ample rational line bundle
on \(Z\), and let \(\theta=c_1(H)\). For every bounded-above family
\((\varphi_i)\) in the genuine \(\operatorname{PSH}(Z,\theta)\) class,
its upper semicontinuous upper envelope \(\sup_i^*\varphi_i\) belongs to
that same class. Equivalently, for every continuous obstacle \(f\),
\(P_\theta(f)\) is continuous and genuinely \(\theta\)-psh. The
normalized genuine PSH space is compact for pointwise convergence on
divisorial valuations.

This lemma, including first membership on the original target and its
arbitrary-cardinality extension, is established background here. The
following proof route specifies the content being used. It is not a
deduction from a smooth quotient assumption, nor from a source headline
alone.

Over a countable field, take a projective birational resolution
\(\pi:\widetilde Z\to Z\) with smooth projective source. Its morphism is
birational and projective; it need not be smooth. Set \(M=\pi^*H\), an
actual semiample nef rational bundle. The first-membership argument
descends every actual \(M\)-psh potential \(\varphi\) to an upper
semicontinuous scalar \(\psi\) on the original \(Z\), with
\(\varphi=\psi\circ\pi\) at \textbf{every} analytic point, including
values \(-\infty\). This all-point fiber statement is part of the
established descent input; equality on a dense open set and upper
semicontinuity alone would not prove it.

The reference in that descent is literal. Choose \(qH\) very ample and
an actual generating system \(s_0,\ldots,s_N\). On the original affine
chart \(U_i=D(s_i)\), use its original algebra generators \(s_j/s_i\)
and the frame \(s_i\). Then \[
 \rho_i=q^{-1}\max_j\log|s_j/s_i|
        =-q^{-1}\log\|s_i\|_{\rm triv}\geq0,
 \qquad u_i=\psi+\rho_i.                                        \tag{2.1}
\] The constant ratio is present. The canonical zero-weight metric is
zero at every analytic point, since all section norms are at most one
and a generating section is a unit at the valuation center. Pulling
(2.1) back gives exactly \(\varphi+\rho_i\circ\pi\) in the same frame.

Here the local conversion of a nef potential uses actual bundles, not
only numerical approximations. Decreasing rational ample-class
approximants \(\varphi_a\) can be put in \(M+\delta A\), with \(A\)
actual very ample, by adding zero in the ample difference class. The
valid pairs \((a,\delta)\) are directed by increasing \(a\) and
decreasing positive rational \(\delta\). In fixed frames the raw weights
are \(\varphi_a+h_M+\delta h_A\) with \(h_A\geq0\), and decrease to
\(\varphi+h_M\). At each valid pair, same-actual-bundle FS approximation
and summable positive rational shifts give decreasing raw FS weights.
Compact directed refinement flattens the two levels into a single
decreasing net. At a \(-\infty\) limit the comparison is with an
arbitrary finite threshold. Every missing original algebra generator and
a constant floor can be appended with finite weights below the old
finite minimum on the compact neighborhood. Thus this operation retains
the original ring and the reference exactly.

On a countable finite-type source, second countability extracts an
increasing sequence of indices with the same pointwise infimum on a
compact seminorm neighborhood, without requiring cofinality. Choose a
smaller compact subset inside its ordinary interior whenever a
surrounding common open is needed. A winning finite logarithmic branch
is nonvanishing and becomes a unit on a sufficiently small principal
neighborhood; the constant floor supplies the constant support.
Vanishing suppressed branches remain suppressed there. Rational shifts
apply to both weights and supports. Restrictions, products, and finite
combinations retain the same relative proper witnesses and original
generators; finite common projective domination is used for finitely
many witnesses. These are the original-reference contact antecedents of
the established countable first-membership theorem.

Its projective consequence puts, for each rational \(c>1\), the chart
function \(\psi-(c-1)\rho_i\), extended by \(-\infty\) outside \(U_i\),
in the actual \(cH\)-psh class. At every point choose a maximal
homogeneous coordinate; on that chart \(\rho_i=0\). Every other chart
contribution is at most \(\psi\). Their finite maximum is therefore
exactly \(\psi\), including at \(-\infty\). Consequently
\(\psi\in\operatorname{PSH}(cH)\) for all \(c>1\). Varying-class
decreasing closure applied to the \textbf{constant scalar family}
\(\psi\) as \(c\downarrow1\) gives its first genuine \(H\) membership.
The individual chart extensions increase and are not the decreasing
family used at this last step. Pullback is now an onto bijection, the
divisorial coordinates in the common function field identify its
topologies, and analytic surjectivity preserves suprema. Smooth-source
nef compactness, transported through this actual identification, gives
target compactness. The general-projective envelope/compactness
equivalence then gives target EP. No target EP was used to obtain first
membership.

For completeness, the arbitrary-cardinality reduction retains the
original possibly nonnormal variety. For a finite rational PL obstacle
\(f\), descend \(Z,H,f\), an ample embedding, and the finite
normalization with its algebra structure, overlap maps, and inverse on a
dense open to a countable algebraically closed subfield \(k_0\).
Normality descends by faithful flatness; the finite normal birational
algebra is the actual normalization. Its universal homeomorphism
descends fpqc locally, so the original model \(Z_0\) is geometrically
unibranch. The countable conclusion applies to the actual \(H_0\).
Restriction \(r_0:Z_k^{\rm an}\to Z_0^{\rm an}\) is surjective: a
radius-one Gauss extension of \(\mathcal H(x)\) on a transcendence
basis, followed by an algebraic norm extension and completion, lifts any
analytic point. Every polynomial uses finitely many variables. The
embedded ground field is still trivially valued, by the dominant-term
test in a monic equation with nonzero constant term.

The pulled envelope \(\psi=P_{H_0}(f_0)\circ r_0\) is continuous and
genuine in the original bundle. Continuous raw Monge--Ampere base-change
compatibility and countable orthogonality give \[
       \int(f-\psi)\,\operatorname{MA}_{\rm raw,H}(\psi)=0.       \tag{2.2}
\] There is no assertion about equality of inverse-image supports. Each
finite competitor \(h\in H^{\rm dom}(H)\) is a literal section formula
on the same actual multiple of \(H\). Descend its finitely many
coefficients to its own larger countable field \(k_2\). A finite-cover
equalizer gives section base change, and faithful flatness descends
global generation. Surjectivity of restriction transfers \(h\leq f\).
The continuous genuine pullback of \(\psi\) is bounded, hence finite
energy, over \(k_2\) without another inverse existence theorem. Raw
pushforward transfers (2.2), and the general-projective domination
principle gives \(h\leq\psi\). For an arbitrary genuine competitor
\(q\leq f\), its decreasing \(H^{\rm dom}(H)\) majorants \(h_a\) satisfy
\(h_a\leq f+\epsilon\) eventually: the decreasing compact sets
\(\{h_a\geq f+\epsilon\}\) have empty intersection. Rational shifts
reduce this to the finite-competitor argument. Letting
\(\epsilon\downarrow0\) proves the exact envelope identity on the
original target. Uniform PL density, the envelope Lipschitz bound, and
genuine uniform closure extend the result to every continuous obstacle.
No arbitrary family is assigned a single countable field.

Primary provenance for this route is the projective theory of
Boucksom--Jonsson {[}BJ, Sections 3--5, 7--10, 13{]}, its same-bundle
regularization, varying-class closure, domination and orthogonality; the
continuous raw base-change construction in Boucksom--Eriksson {[}BE{]};
and the original-reference countable contact and descent route
associated to Wu {[}W26, Section 7{]}. The all-point descent and
repaired first-membership join, together with the discharged
countable-to-arbitrary implication, are the established mathematical
background being retained, rather than an inference from the printed
countable headline alone. In particular, {[}BJ, Theorem 5.20{]} by
itself proves smooth EP; it does not prove singular target EP.

We record the projective consequences used below, with their actual
scopes. For an irreducible projective variety and an ample class, EP
implies normalized compactness by {[}BJ, Theorem 5.11{]}. The proof
bounds every normalized divisorial coordinate, takes a Tychonoff subnet,
and forms the genuine decreasing upper tail envelopes
\(\sup_{j\geq i}^*\varphi_j\). Nonnegligibility of divisorial points
identifies their finite limit; decreasing PSH closure produces an actual
potential. This works for arbitrary nets and uncountable fields. For
every continuous \(\gamma\), {[}BJ, Proposition 4.28{]} makes
\(\varphi\mapsto\sup(\varphi-\gamma)\) continuous. Its proof detects a
PL obstacle on finitely many divisorial points and then uniformly
approximates a continuous obstacle. {[}BJ, Theorem 7.1 and Proposition
7.7{]} give energy upper semicontinuity and decreasing continuity on
general projective varieties. These results require neither smoothness
nor inverse existence.

Two further general-projective inputs will be important. {[}BJ, Theorem
9.20{]} equates finite dual energy with integrability of every actual E1
potential on an irreducible projective variety. {[}BJ, Proposition
9.19{]} says that a divisorially converging genuine PSH net with
eventually bounded \(J\) converges in \(L^1(\lambda)\) for every
finite-dual-energy projective probability \(\lambda\). It does not
require EP. Its signed-integral assertion uniformly approximates
integration on bounded-\(J\) sets by pairings with bounded model MA
probabilities. Those finite divisorial pairings are weakly continuous;
the approximation uses maximizing model outputs and outer
differentiation, not an inverse equation. For the absolute assertion
apply the signed assertion also to \(\max(\varphi_i,\varphi)\), whose
\(J\) is bounded by monotonicity. The exact identity is \[
 |a-b|=2\bigl(\max(a,b)-b\bigr)-(a-b).                           \tag{2.3}
\] The opposite sign in the extracted printed proof is not used. Energy
upper semicontinuity first makes the bounded-\(J\) limit E1.

The relative-volume estimate supporting that differentiability route has
geometrically reduced projective scope over arbitrary complete
non-Archimedean fields {[}BGM, Lemma 3.2{]}. It passes to an
algebraically closed nontrivially valued extension, approximates the
same actual metrics by nef models, takes a reduced-special-fiber model,
and perturbs toward actual ample models. Relatively regular sections and
restriction to effective vertical divisors give a content bound; its
asymptotic coefficient is \(m^p/p!\). Relative volume uses
\(p!m^{-p-1}\), so both signed inequalities have the stated
normalization. Replacement of the direction by its negative gives both
derivative signs. For a trivial ground field, model-function density is
used after the nontrivial extension, with base-change invariance.
{[}BGM, Section 3.2{]} separately assumes envelope continuity; that
extra assumption is not erased. Thus none of these source applications
introduces smoothness, countability, or inverse existence into
Proposition 9.19.

\section{3. The original normal rational quotient satisfies the envelope
lemma}\label{the-original-normal-rational-quotient-satisfies-the-envelope-lemma}

Choose homogeneous algebra generators \(f_1,\ldots,f_N\) of \(R\), of
nonzero weights \(\beta_1,\ldots,\beta_N\). The finite inequalities \[
                 \langle\eta,\beta_i\rangle>0\quad(1\leq i\leq N)
                                                               \tag{3.1}
\] define a nonempty open rational cone, because the prescribed \(\xi\)
belongs to it. A rational vector exists in this open set. Every
occurring nonzero weight is a weight of a monomial in the generators, so
(3.1) supplies all required positivity. Clear denominators, then divide
all supported degrees by their common gcd. We obtain an effective
integral positive grading \[
 R=\bigoplus_{m\geq0}R_m,\qquad R_0=k,\qquad
                     \eta=s\eta_e,\quad s>0.                    \tag{3.2}
\] The notation \(\eta_e\) refers to this effective grading action. If
the original cocharacter lattice has a finite ineffective kernel, it is
the effective quotient action that is used; no lattice primitivity is
silently inferred before removing the supported gcd.

Set \(Y=\operatorname{Proj}R_{\eta_e}\). It is the original integral
projective quotient, of dimension \(p=n-1\). The positive finitely
generated grading gives projectivity and a sufficiently divisible \(d\)
for which \[
                       L_d=\mathcal O_Y(d)                     \tag{3.3}
\] is an actual invertible ample bundle. One elementary description is
to take \(d\) divisible by every generator degree. The powers of the
generators of degree \(d\) define a standard graded subalgebra over
which \(R\) is finite: each original generator satisfies its monic power
equation. On the covering charts their degree-\(d\) powers trivialize
(3.3). This gives the finite map to the standard projective quotient,
its ample bundle, and the dimension formula. Increasing divisibility
gives the usual Veronese construction and, if desired, very ampleness.
All section systems below use positive powers of this same bundle.

Here is normality on the original quotient, without replacing it by a
normalization. For a homogeneous \(f\), the localization \(R_f\) is
normal. If \(z\in\operatorname{Frac}((R_f)_0)\) is integral over
\((R_f)_0\), it is integral over \(R_f\) and hence belongs to \(R_f\).
As a rational degree-zero element it is invariant under the grading
action. The direct graded decomposition, and the infinite ground field,
imply \(z\in(R_f)_0\). Therefore every chart of \(Y\) is a normal
domain. Its local rings at every point are normal domains. Their
reductions and integral closures are the same local rings, and the
residue-field extensions are identities, hence purely inseparable. The
definition of geometric unibranchness applies at every point; this is
exactly {[}Stacks, 0BPZ and 0BQ3{]}.

Consequently \(Y\), over the original field of its original arbitrary
cardinality, and the actual ample rational bundle \(L_d\), satisfy
\textbf{all} the hypotheses of the established projective envelope
lemma. EP holds on \(Y\) itself. No resolution of \(Y\) is substituted
for this application. If further positive tensor scaling is needed, the
literal correspondence
\(\operatorname{PSH}(a\theta)=a\operatorname{PSH}(\theta)\) gives \[
                    P_{a\theta}(f)=aP_\theta(f/a),\quad a>0,    \tag{3.4}
\] and preserves EP. Thus the choice of sufficiently divisible \(d\)
leaves no additional envelope hypothesis.

The normal integral \(Y\) is reduced and of finite type over the perfect
field \(k\), hence geometrically reduced. It is projective, and \(L_d\)
is actual ample. These also verify the general-scheme antecedents of the
relative-volume estimate supporting the projective target-continuity
theorem; no smoothness is being inferred from normality.

The quotient retains the prescribed torus. The natural
\(T\)-linearization of \(L_d\) descends to the residual torus after
twisting by a degree-\(d\) eigencharacter \(\gamma\): on sections of
\(mL_d\) the weights are \(\alpha-m\gamma\), so the effective central
grading acts trivially. Eigen-section frames give degree-zero residual
graded local rings. Taking the Gauss component minimum defines the
continuous residual torification \[
 q:Y^{\rm an}\longrightarrow Y^{\rm an},\qquad B=q(Y^{\rm an}),\qquad
 q^2=q.                                                        \tag{3.5}
\] The frame transitions are eigenunits and make these maps compatible.
\(B\) is compact and closed. The canonical trivial reference is
preserved. The calibrated rational cone correspondence identifies \(B\)
with the original \(K\) by a homeomorphism \(\ell:B\to K\), including
points with kernels. Degree-zero ratios determine the base point, and
the radial Gauss coordinate supplies the cone extension. These are the
original quotient data and references of the continuous cone
correspondence.

It is important to retain every calibration factor: \[
 \begin{split}
 V_\eta&=s^{-n}d^{1-n}(L_d^p),\\
 \omega_\eta&=L_d/(sd),\\
 (\omega_\eta^p)&=sV_\eta,\\
 E_\eta(u)&=(sd)^{-1}E_{L_d}(\chi),\\
 \chi&=sd\,u\circ\ell\circ q.
 \end{split}                                                   \tag{3.6}
\] At the continuous level these are the calibrated Gauss/finite-map
identities: the degree-\(d\) finite map and the \(n\) factors \(1/(sd)\)
give the radial normalization. The projective mixed probability is raw
projective curvature divided by \((L_d^p)\). That denominator is not
\(V_\eta\).

\section{4. Torification and both genuine unbounded rational
correspondences}\label{torification-and-both-genuine-unbounded-rational-correspondences}

We now verify the extension of the continuous correspondence, which is
the point at which an abstract completion would be inadequate.

\textbf{Lemma 4.1.} For every genuine projective PSH potential \(\chi\)
on \(L_d\), \(\chi^T=\chi\circ q\) is genuine PSH and
\(\chi^T\geq\chi\). It preserves normalization, and \[
                  E_{L_d}(\chi^T)\geq E_{L_d}(\chi).             \tag{4.1}
\] It therefore takes actual projective E1 into actual invariant E1.

For an FS formula on the actual \(mL_d\), write each section as a finite
sum of residual eigen-sections. Retain the additive coefficient of the
section for each of its components. The components still globally
generate: simultaneous vanishing of all components would force
simultaneous vanishing of all original sections. In an eigenframe the
torified norm is the maximum of the component norms. The ultrametric
inequality bounds the norm of the original sum by that maximum. Thus the
resulting actual FS formula is \(\chi\circ q\geq\chi\), on the same
\(mL_d\) and reference. Uniform approximation extends this to genuine
CPSH. If \(\chi_a\downarrow\chi\) is a genuine decreasing CPSH defining
net, then \(\chi_a\circ q\downarrow\chi\circ q\) at every point,
including values \(-\infty\). The generic trivial valuation is fixed and
its value is finite, so this is a nontrivial genuine PSH limit.
Monotonicity of the actual energy proves (4.1). Normalization at the
trivial valuation is unchanged.

\textbf{Lemma 4.2.} The continuous rational correspondence extends to an
onto order-preserving correspondence between actual cone PSH at \(\eta\)
and actual invariant projective PSH on \(L_d\). Formula (3.6) holds also
for unbounded potentials, including energy \(-\infty\). In particular it
is onto actual E1 in both directions.

Indeed, for actual cone \(u\) choose raw decreasing defining potentials
\(u_a\in\operatorname{CPSH}(\eta)\). Their invariant projective lifts
\(\chi_a=sd\,u_a\circ\ell\circ q\) decrease at every analytic point to
\(\chi=sd\,u\circ\ell\circ q\). Its finite value at the trivial
valuation excludes the identically \(-\infty\) limit. Hence \(\chi\) is
actual PSH. Conversely take an invariant actual projective \(\chi\). Its
genuine decreasing CPSH defining net can be torified by Lemma 4.1. The
resulting continuous invariant net decreases pointwise to
\(\chi\circ q=\chi\). Its continuous cone lifts decrease on \(K\) and on
the punctured cone through their actual equivariant extensions. Every
point has the same limiting formula, even when the value is \(-\infty\).
The finite trivial-valuation value again excludes a trivial limit. This
gives an actual cone PSH potential. The two formulas are inverse because
\(q\) is a retraction and \(\ell\) is the calibrated homeomorphism.
Decreasing continuity of both energies extends (3.6), including
\(-\infty\). It proves the E1 assertion. Nested defining nets can be
flattened by compact directed refinement {[}BJ, Lemma 4.6{]}; no
sequential regularization is imposed on an uncountable space. For these
passages we use raw decreasing nets. Subtracting their varying
normalization constants need not preserve decreasingness.

There is also the required actual weighted operator correspondence. Let
\(m_{L_d}(\chi_1,\ldots,\chi_p)\) denote the projective mixed
probability, with denominator \((L_d^p)\), for the invariant E1 lifts of
a cone E1 tuple. Then \[
 \begin{split}
 \ell_*m_{L_d}(\chi_1,\ldots,\chi_p)
     &=\operatorname{MA}_\eta(u_1,\ldots,u_p),\\
 \int_Ku_0\,d\operatorname{MA}_\eta(u_1,\ldots,u_p)
     &=(sd)^{-1}\int_B\chi_0\,dm_{L_d}(\chi_1,\ldots,\chi_p).
 \end{split}                                                   \tag{4.2}
\] All projective measures here are supported on \(B\). At the
continuous stages these are the calibrated rational identities. Choose
simultaneous raw decreasing CPSH defining nets for the actual E1 inputs.
Fixed-\(\eta\) decreasing continuity of the actual cone mixed operator
and the general-projective mixed operator passes the measures to their
limits. Every E1 test is integrable against these mixed E1 measures. The
projective weighted theorem {[}BJ, Proposition 7.21{]} and the
corresponding fixed-field cone theorem pass the weighted integrals to
the limit.

To specify the sign at this step, for projective ample classes
\(\theta_i\) the raw weighted integral is \[
 \begin{split}
 \int\chi_0\,\bigwedge_{i=1}^p(\theta_i+dd^c\chi_i)
  ={}&(\theta_0,\chi_0)\cdot(\theta_1,\chi_1)\cdots(\theta_p,\chi_p)\\
     &-(\theta_0,0)\cdot(\theta_1,\chi_1)\cdots(\theta_p,\chi_p).
 \end{split}                                                   \tag{4.3}
\] It is the test-containing pairing \textbf{minus} the zero-test
pairing, both at the finite stages and at the limit. In the proof of
weighted convergence, convergence of these masses and positivity pass
the weighted measures, not just bounded test integrals. Closedness of
\(B\) and \(K\) preserves support under weak limits. The mass remains
one and the weight factor in (4.2) remains \(1/(sd)\). No operator has
been transported between different Reeb fields in this argument.

\section{5. Intrinsic target integrability and the exact midpoint
bound}\label{intrinsic-target-integrability-and-the-exact-midpoint-bound}

We give the finite-target and midpoint arguments to make clear that they
do not contain the missing realization step.

For a bounded-above decreasing net of upper semicontinuous functions
\(h_a\downarrow h\) on compact Hausdorff \(K\), Radon regularity gives
\[
                        \int h_a\,d\mu\downarrow\int h\,d\mu.    \tag{5.1}
\] For bounded functions use continuous majorants and regularity;
truncation below gives the extended case. If \(g\geq h\) is continuous,
the increasing open sets \(\{h_a<g+\epsilon\}\) cover \(K\). Compactness
and directedness give an entire tail with \(h_a\leq g+\epsilon\). Taking
the infimum over continuous majorants proves (5.1). This is a net
statement, not a substitution of the sequential monotone convergence
theorem.

Uniform FS approximation first equates (1.5) with the CPSH supremum. If
\(L_\xi(\mu)<\infty\) and \(u_a\downarrow u\in\mathcal E^1(\xi)\) is an
actual defining net, then \[
                  \int u_a\,d\mu\geq E_\xi(u_a)-L_\xi(\mu).
\] Decreasing energy continuity and (5.1) show that every E1 \(u\) is
integrable. They also show \[
 L_\xi(\mu)=\sup_{u\in\mathcal E^1(\xi)}
                    \left(E_\xi(u)-\int_Ku\,d\mu\right).        \tag{5.2}
\] If the bounded-test supremum is infinite, the E1 supremum is already
infinite, so (5.2) also holds as an extended functional.

Conversely suppose that every actual E1 potential is \(\mu\)-integrable.
Write \(\rho=\rho_\xi\). For normalized \(u\) let \[
 b_l=-\int_Ku\,d\operatorname{MA}_\xi(u^{[l]},\rho^{[p-l]}),
 \quad 0\leq l\leq p.
\] These are finite, nonnegative, \(b_0=0\), and
\(J_\xi(u)=n^{-1}\sum_l b_l\). Along the actual segment
\(u_t=tu+(1-t)\rho\) the positive mixed expansion gives \[
 \begin{split}
 J_\xi(u_t)
 &=\frac1n\sum_{l=1}^p t\sum_{j=l}^p
                    \binom{j}{l}t^l(1-t)^{j-l}b_l\\
 &\leq 2^nt^2J_\xi(u),\qquad 0\leq t\leq1.
 \end{split}                                                   \tag{5.3}
\] Use \(t^{l+1}\leq t^2\) and
\(\sum_{j=l}^p\binom{j}{l}=\binom{n}{l+1}\leq2^n\). The formula is
initially continuous and extends by actual fixed-field weighted
decreasing continuity.

There is a finite constant \(C\) bounding \(\int(-u)\,d\mu\) on
normalized \(J_\xi(u)\leq1\). Otherwise choose such \(u_j\) with
integrals at least \(2^j\). The actual convex sums with weights
\(2^{-j}\) and the remaining weight on \(\rho\) decrease, have value
zero at the trivial valuation and energy at least \(-1\). Their
nontrivial decreasing PSH limit is actual E1, but has infinite negative
integral, a contradiction. Apply this bound to \(u_t\) with
\(t=(2^nJ_\xi(u))^{-1/2}\) when \(J_\xi(u)\geq2^{-n}\), and directly
otherwise. It follows that \[
 \int(-u)\,d\mu\leq C+C\sqrt{2^nJ_\xi(u)},\qquad
 L_\xi(\mu)\leq C+2^{n-2}C^2<\infty.                          \tag{5.4}
\] This proves the intrinsic criterion without inverse existence.

For \(p\geq1\), let \(z=u-v\) for continuous same-field potentials and
put \[
 B_j=-V_\xi^{-1}\int_K z\,dd^cz\,(dd^cu)^j(dd^cv)^{p-1-j}
                         dd^c\rho_\xi^+\geq0,\quad 0\leq j<p.
\] The existing Hodge inequality gives positivity. Integration by parts
and the full difference-of-powers telescope, with \(j=0,\ldots,p-1\),
give \[
 I_\xi(u,v)=\sum_{j=0}^{p-1}B_j,\qquad
 -\frac{d^2}{dt^2}E_\xi(v+tz)
   =p\sum_{j=0}^{p-1}\binom{p-1}{j}t^j(1-t)^{p-1-j}B_j.
\] The midpoint Green kernel is \(\tfrac12\min(t,1-t)\). On
\([1/4,3/4]\) it is at least \(1/8\), each monomial is at least
\(4^{-(p-1)}\), each binomial coefficient is at least one, and the
interval has length \(1/2\). Therefore \[
 E_\xi((u+v)/2)-\frac{E_\xi(u)+E_\xi(v)}2
               \geq a_pI_\xi(u,v),\qquad a_p=p/4^{p+1}.          \tag{5.5}
\] Every expression is a finite combination of the genuine weighted
mixed pairings, so raw simultaneous decreasing approximation proves
(5.5) also for actual E1. With
\(d_u=L_\xi(\mu)-E_\xi(u)+\int u\,d\mu\geq0\), the deficit of the
midpoint is nonnegative, whence \[
 I_\xi(u,v)\leq\frac{d_u+d_v}{2a_p}.                            \tag{5.6}
\] The reference has deficit \(L_\xi(\mu)\). Its tangent deficit against
a normalized \(u\) is \(J_\xi(u)\), one of the two nonnegative deficits
adding to \(I_\xi(u,\rho_\xi)\). Thus \[
 J_\xi(u)\leq I_\xi(u,\rho_\xi)
                \leq\frac{d_u+L_\xi(\mu)}{2a_p}.               \tag{5.7}
\] This supplies genuine bounded-energy maximizing sequences; by itself
it does not supply an actual limit.

\section{6. A topology on actual potentials and compactness at the
quotient}\label{a-topology-on-actual-potentials-and-compactness-at-the-quotient}

Let \(\mathcal P_\theta\) be the normalized actual PSH class at a Reeb
vector \(\theta\). For every actual continuous test \(\psi\) at that
field, set \[
              S_\psi(u)=\sup_K(u-\psi)\in\mathbb R,
\] and let \(\tau_\theta\) be the topology generated by these maps. The
supremum is finite, since \(u\leq0\), its trivial-valuation value is
zero, and \(\psi\) is bounded on \(K\).

This topology is Hausdorff. If \(u(x)>v(x)\), choose from a genuine
defining net \(\psi_a\downarrow v\) a test with \(\psi_a(x)<u(x)\), even
when \(v(x)=-\infty\). Then \(S_{\psi_a}(u)>0\) while
\(S_{\psi_a}(v)\leq0\). Exchange the functions for the opposite
inequality. Restriction determines the whole potential, so this
separates the actual functions. Energy is upper semicontinuous in
\(\tau_\theta\): if \(E_\theta(u)<c\), choose an actual continuous
majorant \(\psi\geq u\) with \(E_\theta(\psi)<c-\epsilon\). The
neighborhood \(S_\psi(v)<\epsilon\) gives \(v\leq\psi+\epsilon\) and
\(E_\theta(v)<c\). The argument includes energy \(-\infty\).

By Sections 2 and 3 the actual normalized projective PSH space for
\(L_d\) is compact in its divisorial topology. Define \[
 A(\chi)=(sd)^{-1}\chi^T\circ\ell^{-1}\quad\hbox{on }K.         \tag{6.1}
\] Lemmas 4.1 and 4.2 prove actual membership and surjectivity of \(A\).
For a cone continuous test \(\psi\), its invariant projective lift is
\(\Psi=sd\,\psi\circ\ell\circ q\). The exact global identity is \[
 S_\psi(A\chi)=(sd)^{-1}\sup_{Y^{\rm an}}(\chi-\Psi).           \tag{6.2}
\] Indeed \(\Psi\) is invariant, \(q\) is a retraction, and
\(\chi^T\geq\chi\). The supremum of \(\chi^T-\Psi\) equals the supremum
of \(\chi-\Psi\) on \(B\), which is at most the full supremum.
Domination gives the reverse inequality. {[}BJ, Proposition 4.28{]}
makes (6.2) continuous. Thus \(A\) is a continuous surjection from a
compact space onto the Hausdorff space \((\mathcal P_\eta,\tau_\eta)\),
which is consequently compact. No closedness assertion about an
invariant weak subset is needed.

\textbf{Proposition 6.1.} If \(u_i\to u\) in \(\tau_\eta\) and
\(J_\eta(u_i)\) is eventually bounded, then \(u\) is actual E1 and \[
                  u_i\longrightarrow u\quad\hbox{in }L^1(\lambda)
                                                               \tag{6.3}
\] for every cone probability \(\lambda\) with finite intrinsic dual
energy. The assertion concerns the whole original net.

Energy upper semicontinuity first gives \(E_\eta(u)\geq-C\) on the
bounded tail, so \(u\) is actual E1. Push \(\lambda\) through
\(\ell^{-1}\) to \(\lambda_Y\), viewed as a projective probability
supported on \(B\). It integrates \textbf{every} projective E1 potential
\(\chi\): Lemma 4.1 makes \(\chi^T\) invariant E1; Lemma 4.2 gives its
actual cone E1 restriction; Section 5 gives its
\(\lambda\)-integrability; and \(\chi=\chi^T\) on \(B\). Thus {[}BJ,
Theorem 9.20{]} applies on the irreducible normal \(Y\) and makes
\(\lambda_Y\) a finite-projective-dual-energy probability. Knowing only
integrability of invariant projective potentials without this
torification argument would not verify that antecedent.

Suppose (6.3) fails. Passing to a cofinal subnet, its \(L^1\) distances
can be kept above some \(\epsilon>0\). Its invariant projective lifts
\(\chi_i=sd\,u_i\circ\ell\circ q\) are normalized actual E1 and satisfy
\[
                         J_{L_d}(\chi_i)=sdJ_\eta(u_i).
\] Compactness gives a further projective weakly converging subnet
\(\chi_i\to\chi\), with actual normalized PSH limit. Energy upper
semicontinuity makes \(\chi\) actual E1. Continuity of \(A\) and
Hausdorffness identify \(A\chi=u\). No invariance of the weak limit is
assumed. {[}BJ, Proposition 9.19{]}, with the already verified target
\(\lambda_Y\), gives \(\chi_i\to\chi\) in \(L^1(\lambda_Y)\). On \(B\),
\(\chi=\chi^T\) and \(A\chi=u\), so this distance is exactly \(sd\)
times the alleged cone distance. That is a contradiction. Failure of
convergence always admits such a cofinal subnet; hence the entire
original net converges as claimed.

\section{7. Actual root homeomorphisms and the transfer to a fixed
irrational
field}\label{actual-root-homeomorphisms-and-the-transfer-to-a-fixed-irrational-field}

We verify all transfer antecedents, rather than transporting an operator
between Reeb fields. Fix any prescribed rational or irrational \(\xi\)
and the rational \(\eta\) of Section 3. For the original homogeneous
algebra generators set \[
 a=\min_i\frac{\langle\xi,\beta_i\rangle}
                    {\langle\eta,\beta_i\rangle}>0,\quad
 b=\max_i\frac{\langle\xi,\beta_i\rangle}
                    {\langle\eta,\beta_i\rangle}<\infty,\quad
 \delta=\max_i\left|\log\frac{\langle\xi,\beta_i\rangle}
                            {\langle\eta,\beta_i\rangle}\right|.
                                                               \tag{7.1}
\] The existing actual cone comparison calculus gives, for an
\(\eta\)-potential and \(t\geq0\), \[
 u(v)-bt\leq u(v+t\xi)\leq u(v)-at.                            \tag{7.2}
\] Reverse the orbit to obtain the bounds for negative \(t\). A finite
orbit is Lipschitz and strictly decreasing; an infinite-valued orbit is
identically \(-\infty\). The zero parameter of a finite orbit is the
raw root \(\mathcal R_{\eta\xi}u\); on an identically \(-\infty\)
orbit, the root has the singular value supplied by the decreasing-limit
extension. The actual root maps preserve order, CPSH,
normalization, and decreasing PSH nets, and are inverse bijections. On
normalized actual E1 they are onto and satisfy \[
 e^{-(2n+1)\delta}J_\eta(u)
 \leq J_\xi(\mathcal R_{\eta\xi}u)
 \leq e^{(2n+1)\delta}J_\eta(u),\qquad
       u/a\leq\mathcal R_{\eta\xi}u\leq u/b\leq0\ \hbox{on }K.
                                                               \tag{7.3}
\] These are comparisons of genuine functions and energies. By the
intrinsic integrability criterion of Section 5 and the onto inverse, the
finite-energy target class on the fixed \(K\) is the same for all Reeb
fields.

For completeness the actual root bijection is a homeomorphism for the
supremum-test topologies. If \(\psi\in\operatorname{CPSH}(\xi)\), set
\(g_t=\mathcal R_{\xi\eta}(\psi+t)\). Raw global order equivalence gives
\[
 S_\psi(\mathcal R_{\eta\xi}u)\leq t
              \quad\Longleftrightarrow\quad S_{g_t}(u)\leq0.    \tag{7.4}
\] It includes unbounded values, since order on restrictions is order
after the equivariant extensions. Strict negative gaps are preserved in
both directions: by (7.2), adding a positive constant \(c\) changes
every finite root by an amount between \(c/b\) and \(c/a\), with the
same order bounds at \(-\infty\). The inverse ratios are \(1/b\) and
\(1/a\). Hence (7.4) also holds with \(<\) on both sides, and
complements give \(>\) on both sides. All strict sublevel and superlevel
preimages are open. Applying the argument to the inverse proves the
homeomorphism. It transports actual normalized compactness from Section
6 to \((\mathcal P_\xi,\tau_\xi)\).

We next check target continuity. For each real \(r\) define on the
actual carrier \[
 T_r(x)=s_\eta(x+r\xi),\qquad
       s_\eta(v)=v+\rho_\eta(v)\eta.                            \tag{7.5}
\] This map is continuous; the torified action and the reference
retraction are continuous. No inverse or group law for \(T_r\) is
required. Equivariance and (7.2), applied also to the reference, give
for finite \(u(x)\) \[
 \begin{split}
 u(T_rx)&=u(x+r\xi)-\rho_\eta(x+r\xi),\\
                         |u(T_rx)-u(x)|&\leq2b|r|.
 \end{split}                                                   \tag{7.6}
\] The \(-\infty\) locus is preserved: the entire orbit remains
\(-\infty\), and subtracting the finite reference does not change it. If
\(\mu\) is a finite-energy probability, every actual \(\eta\)-E1
potential is therefore integrable against \(\mu_r=(T_r)_*\mu\). Section
5 makes \(\mu_r\) a finite-energy target. Continuous pushforward on
compact Hausdorff spaces preserves positivity, Radon regularity and mass
one.

Here is the full finite-grid and tail argument. Suppose normalized
actual \(u_i,u\in\mathcal E^1(\eta)\) satisfy \[
 u_i\to u\text{ in }L^1(\mu),\qquad
 u_i\to u\text{ in }L^1(\mu_r)\quad\text{for every rational }r\leq0.
                                                               \tag{7.7}
\] Put \(R_i=\mathcal R_{\eta\xi}u_i\), \(R=\mathcal R_{\eta\xi}u\), and
\[
 G_i(r,x)=u_i(x+r\xi)=u_i(T_rx)+\rho_\eta(x+r\xi),\qquad
 G(r,x)=u(x+r\xi).
\] The finite orbit functions have slopes between \(-b\) and \(-a\).
Null sets may depend on the net index: each inequality at a fixed index
involves only finitely many functions, so it needs no common null set
for an uncountable net. The roots are nonpositive, with \[
                    0\leq-R_i\leq-u_i/a,\qquad0\leq-R\leq-u/a. \tag{7.8}
\] Fix positive rational \(M\) and a rational mesh \(h=M/N\) on
\([-M,0]\), with grid \(r_j=-M+jh\). The Lipschitz bound gives \[
 \sup_{-M\leq r\leq0}|G_i(r,x)-G(r,x)|
             \leq\max_j|G_i(r_j,x)-G(r_j,x)|+2bh.
\] For two strictly decreasing functions whose downward slope is at
least \(a\), their zero parameters clipped to \([-M,0]\) differ by at
most their supremum discrepancy on that interval divided by \(a\). For
an interior zero use the slope bound at the other zero. If a clipped
zero is \(-M\), its value at \(-M\) is nonpositive, and the same
estimate at \(-M\) bounds the distance to the other interior zero. The
zero endpoint is treated in the same way. Thus with
\(R_i^M=\max(R_i,-M)\), \[
 \|R_i^M-R^M\|_{L^1(\mu)}
 \leq a^{-1}\sum_{j=0}^N\|u_i-u\|_{L^1(\mu_{r_j})}+2bh/a.       \tag{7.9}
\] The reference terms cancel. First take the original net limit with
the finite grid fixed; then refine the mesh. This proves clipped \(L^1\)
convergence for each \(M\).

The tails, not just the norms, are controlled by \[
 \|R_i-R_i^M\|_{L^1(\mu)}
                \leq a^{-1}\int_K(-u_i-aM)_+\,d\mu.            \tag{7.10}
\] For fixed \(M\) the integral tends to the one with \(u\), since this
scalar positive-part function is 1-Lipschitz and (7.7) holds at zero.
Combining (7.9) and (7.10) gives \[
 \limsup_i\|R_i-R\|_{L^1(\mu)}
                     \leq\frac2a\int_K(-u-aM)_+\,d\mu.        \tag{7.11}
\] Send \(M\to\infty\) last. Integrability of the fixed actual \(u\)
proves \(R_i\to R\) in \(L^1(\mu)\).

Apply this argument to a \(\tau_\xi\) convergent actual normalized E1
net \(w_i\to w\) with eventually bounded \(J_\xi\). Energy upper
semicontinuity makes \(w\) actual E1. The inverse roots \(u_i,u\) are
actual E1, \(\tau_\eta\) convergent, and have bounded \(J_\eta\) by
(7.3). Every \(\mu_r\) is finite energy by (7.6), so Proposition 6.1
supplies all of (7.7). Equations (7.9)--(7.11) prove convergence of the
entire original net in \(L^1(\mu)\). This verifies the
root-homeomorphism, energy comparison, orbit finite-energy, and
grid/tail antecedents for every fixed irrational \(\xi\). Throughout,
the MA operator remains at its own fixed field.

\section{8. Actual attainment and identification of the genuine
fixed-field
operator}\label{actual-attainment-and-identification-of-the-genuine-fixed-field-operator}

Assume \(p\geq1\) and \(L_\xi(\mu)<\infty\). By the real supremum (1.5)
choose actual \(h_j\in H(\xi)\) with deficits tending to zero, and
normalize them by constants. Equation (5.7) gives eventually bounded
\(J_\xi(h_j)\). Actual normalized compactness from Section 7 yields a
subnet converging in \(\tau_\xi\) to an actual normalized PSH potential
\(w\). Energy upper semicontinuity makes it actual E1, and the whole-net
target-continuity result, applied to this subnet, gives target \(L^1\)
convergence. Therefore \[
 E_\xi(w)-\int_Kw\,d\mu
 \geq\limsup_j\left(E_\xi(h_j)-\int_Kh_j\,d\mu\right)
                      =L_\xi(\mu).                            \tag{8.1}
\] The opposite inequality is (5.2). Hence \(w\) actually attains the
supremum. It is not obtained by completing an \(I\)-Cauchy class.

We identify its \textbf{actual} fixed-\(\xi\) operator. For continuous
obstacles define the genuine outer energy \[
 F_\xi(g)=\sup\{E_\xi(v):v\in\operatorname{CPSH}(\xi), v|_K\leq g\},
 \qquad
 F_\xi^\downarrow(h)=\inf_{g\in C(K),g\geq h}F_\xi(g).
                                                               \tag{8.2}
\] The existing cone calculus makes the latter agree with energy on
actual PSH and be continuous along decreasing obstacles. For each
continuous \(g\), feasibility gives
\(F_\xi(g)-\int g\,d\mu\leq L_\xi(\mu)\). A bounded-above integrable
upper semicontinuous \(h\) is the decreasing net limit of its continuous
majorants on compact \(K\); minima make the family directed. Radon
regularity and (5.1) then give the genuine outer-obstacle inequality \[
                F_\xi^\downarrow(h)-\int_Kh\,d\mu\leq L_\xi(\mu).
                                                               \tag{8.3}
\]

The cone calculus supplies a uniformly dense continuous direction
lattice \(\mathcal D\subset C(K)\) and, for each fixed
\(f\in\mathcal D\), both signs of the fixed-field outer derivative \[
 F_\xi^\downarrow(w+tf)
      =E_\xi(w)+t\int_K f\,d\operatorname{MA}_\xi(w)+O_f(t^2).
                                                               \tag{8.4}
\] Its operator is exactly the genuine probability (1.2). The lattice
can be constructed on one rational ray from rational spans of
differences of FS functions, allowing larger positive equivariant
slopes. Products clear exponents, maxima concatenate lists, and primary
reference padding matches slopes. The identity
\(\max(a-b,c-d)=\max(a+d,c+b)-(b+d)\) gives lattice closure. Homogeneous
functions separate distinct torified points, including finite/infinite
evaluations, and constants are included. Lattice Stone--Weierstrass
gives uniform density. The established Reeb comparison and calibrated
relative-volume estimate furnish (8.4) at every fixed \(\xi\); it is an
outer derivative and needs no enveloped inverse assertion.

Apply (8.3) to \(h=w+tf\), which is bounded above, upper semicontinuous
and \(\mu\)-integrable. Since (8.1) has zero deficit, (8.4) gives \[
 t\left(\int_K f\,d\operatorname{MA}_\xi(w)-\int_K f\,d\mu\right)
                                                        \leq O_f(t^2).
\] Divide by positive \(t\) and let \(t\downarrow0\), then use negative
\(t\) and let \(t\uparrow0\). The two inequalities force equality of
both integrals. Uniform density and the two Radon probabilities of mass
one identify them on all continuous tests. Hence \[
                          \operatorname{MA}_\xi(w)=\mu.        \tag{8.5}
\] No root pushforward of MA and no abstract assigned inverse measure
enters this identification.

Necessity uses the actual tangent inequality, extended by weighted
decreasing continuity: \[
 E_\xi(v)\leq E_\xi(u)+\int_K(v-u)\,d\operatorname{MA}_\xi(u),
 \qquad u,v\in\mathcal E^1(\xi).
\] All E1 tests are integrable against genuine E1 outputs, so \[
 L_\xi(\operatorname{MA}_\xi(u))
                 =E_\xi(u)-\int_Ku\,d\operatorname{MA}_\xi(u)<\infty.
                                                               \tag{8.6}
\] Support, positivity and mass one follow from the genuine operator.
The existing bi-Lipschitz embedding onto its image gives uniqueness
modulo constants. It is used for uniqueness after existence, not as an
inference of surjectivity.

Finally let \((u_j)\) be \textbf{any} normalized actual E1 maximizing
sequence, not only the FS sequence chosen for attainment. Apply (5.6) to
\(u_j\) and the actual zero-deficit maximizer \(w\). This proves exactly
(1.7). In addition (5.7) bounds \(J_\xi(u_j)\). Every compact convergent
subnet has, by energy upper semicontinuity and target \(L^1\)
continuity, an actual zero-deficit limit. Equation (8.5) and uniqueness
identify that limit with \(w\). Compactness and Hausdorffness therefore
give convergence of the whole sequence in \(\tau_\xi\), and Section 7
gives target \(L^1\) convergence. The deficit identity then gives
\(E_\xi(u_j)\to E_\xi(w)\). These conclusions realize the entire
maximizing sequence in the genuine class.

\section{9. The literal one-point
case}\label{the-literal-one-point-case}

If \(n=1\), the original effective normal positive grading has
\(R=k[z]\). One way to see the elementary grading conclusion is that the
degree-zero graded fraction field is the function field of the
zero-dimensional integral Proj, hence \(k\). Each nonzero homogeneous
piece is therefore at most one-dimensional. In a Laurent homogeneous
coordinate its supported degrees form a finitely generated positive
semigroup of gcd one. Normality makes that semigroup saturated: if a
positive multiple of a degree occurs, the corresponding Laurent monomial
is integral and must be in the ring. The semigroup is thus
\(\mathbb N\). This gives \(k[z]\) with the effective grading,
irrespective of the original ineffective torus kernel.

The quotient is the original point \(\operatorname{Spec}k\), whose
normalized PSH class consists of zero and whose EP is immediate. The
calibrated cone carrier is \(K=\{v_{\rm triv}\}\). Every potential
restricts to a scalar \(c\), its energy is \(E(c)=c\), and the empty
horizontal product in (1.2) has the calibrated radial unit mass. The
sole probability \(\delta_{v_{\rm triv}}\) has dual energy zero and is
realized. Constants give uniqueness modulo constants, and every
normalized maximizing sequence is already the zero restriction. There is
no positive-dimensional Hodge or midpoint estimate in this case.

\section{10. Evidence, dependencies, and
reproducibility}\label{evidence-dependencies-and-reproducibility}

The conclusion is for the unchanged original domain. No rationality,
smoothness, field-countability, resolution-EP, or additional quotient-EP
condition remains. The sole newly applied projective geometry is checked
on the original normal quotient in Section 3. The unbounded
correspondences, projective target integrability, whole-net \(L^1\)
argument, irrational transfer, and actual variational attainment are
proved in Sections 4--8. The established background consists of the
exact original cone calculus and the projective envelope lemma, with the
proof route and scope stated in Sections 1--2. Historical review labels
are not proofs of the new application.

For the direct general-projective applications, the pinned primary is
arXiv:1801.08229v4 {[}BJ{]}. Its PDF SHA256 is
\texttt{8d5975337f4ac88a475929ea6a436ec4a15d08715d5582bc0ba0eb3e07152c52};
the supplied literal extract SHA256 is
\texttt{659bb8d2e5335781041f03cdb168bac70126487bc0db394f3dac8cf1e9f076d0}.
Literal statements and pertinent proof bodies for the same-bundle
regularization, supremum continuity, compactness, energy, mixed
measures, weighted convergence, and target \(L^1\)/integrability are
retained locally. The general relative-volume proof is pinned to
arXiv:2004.03847v1 {[}BGM{]}, and normal/geometrically-unibranch
definitions and proofs to the actual Stacks HTTP primary pages. The
absolute-value and weighted-pairing signs are corrected algebraically in
(2.3) and (4.3). For an auxiliary projective PL interpolation use
\(\omega+(1-t)D+tH\) before rational approximation; the printed
\(\omega+D+tH\) is not that interpolation.

\texttt{source-provenance.json} retains versions, exact source and
excerpt hashes, and the literal pertinent bodies. Source coordinates are
inclusive one-based LF indices after UTF-8 byte decoding and
\texttt{split(\textquotesingle{}\textbackslash{}n\textquotesingle{})}.
For the live Stacks rendering, its original zero-based displayed
coordinates are also recorded, separately from the LF coordinates of the
locally preserved text. \texttt{external-premises.json} has nonnegative
premise indices and actual HTTP(S) primary passages covering all
nineteen external premises, including the published inputs in the
recalled established envelope proof route. Task-granted mathematical
background has no supplied native support handle; it is not assigned a
fictitious old note identifier. No full primary paper or every
recursively cited proof is claimed read.

Run \texttt{python3\ rerun.py} in this directory. It checks exact
original input, task and assignment preservation, primary hashes and
literal LF units, finite rational diagnostic inequalities, existing-tool
document generation, and two pdfTeX passes.
\texttt{input-original.json}, \texttt{task-original.json}, and
\texttt{assignment-original.txt} preserve the supplied input. Build
outputs and environment versions are retained in the directory.
\texttt{status.json} has the five required fields;
\texttt{canonical-notes.json} registers fresh local backward note
dependencies. Finite diagnostics are observations on their declared
finite domains, not proofs of the infinite-dimensional argument.
Successful compilation, reading records and hashes are not mathematical
acceptance.

No additional model, worker, installation, private transcript, unrelated
campaign artifact, source-cache mutation, or external message was used.
No resource-heavy computation was needed. No supervised compute facility
was exposed, and none was required or launched.

\section{References}\label{references}

{[}BJ{]} S. Boucksom and M. Jonsson, \emph{Global pluripotential theory
over a trivially valued field}, arXiv:1801.08229v4,
\url{https://arxiv.org/pdf/1801.08229v4}.

{[}BE{]} S. Boucksom and D. Eriksson, \emph{Spaces of norms, determinant
of cohomology and Fekete points in non-Archimedean geometry},
arXiv:1805.01016v3, \url{https://arxiv.org/pdf/1805.01016v3}.

{[}BGM{]} S. Boucksom, W. Gubler and F. Martin, \emph{Differentiability
of relative volumes over an arbitrary non-Archimedean field},
arXiv:2004.03847v1, \url{https://arxiv.org/pdf/2004.03847v1}.

{[}W24{]} Y. Wu, \emph{Some non-Archimedean pluripotential theory on
polarized affine cones}, arXiv:2406.13684v1,
\url{https://arxiv.org/pdf/2406.13684v1}.

{[}W26{]} L. Wu, \emph{On the local theory of non-Archimedean psh
functions}, arXiv:2609.27610v1, Section 7,
\url{https://arxiv.org/pdf/2609.27610v1}. The repaired
original-reference first-membership and cardinality route used here is
the established background specified in Section 2; a printed headline
alone is not the premise.

{[}Stacks{]} The Stacks Project, Definition 15.108.1 and Lemma 28.16.2,
\url{https://stacks.math.columbia.edu/tag/0BPZ},
\url{https://stacks.math.columbia.edu/tag/0BQ3}.

\end{document}
